% ===========================================================================
% SOURCES: docs/claude_project/11_REPO_MAP.md
%          scripts/README.md, scripts/run_*.sh, scripts/*.tcl
%          firmware/README.md, firmware/Makefile
%
% KEEP THIS CHAPTER SHORT - 4 to 6 pages. It is the chapter most likely to
% bloat into product descriptions nobody reads. The test for every sentence:
% does a reader need it to reproduce or to judge a result in chapter 10?
% ===========================================================================
\chapter{Experimental Setup and Development Tools}
\label{chap:setup}

\section{Experimental Setup}
\label{sec:setup}

\subsection{The ZCU104 Evaluation Board}
\label{sec:zcu104}
% [xilinx_ug1267]. The part is XCZU7EV-ffvc1156-2-e. Give the resource table
% once, here, and refer back to it from every utilisation result:
%   DSP48E2 1728 | CLB LUT 230400 | CLB FF 460800 | BRAM 312 | URAM 96
% Section~\ref{sec:ps-pl} covers the architecture; this covers the board.

\subsection{Host Development System}
\label{sec:host}
% CPU, RAM, OS. This matters for exactly one number - the CPU baseline in
% Section~\ref{sec:comparison} - so state it precisely and move on.

\section{Development Tools}
\label{sec:tools}

\subsection{Vivado Design Suite 2022.2}
\label{sec:vivado}
% Version 2022.2.2 exactly. Note that synthesis and implementation are driven
% from scripts (scripts/synth.tcl, scripts/impl.tcl, run via run_synth.sh and
% run_impl.sh), not from the GUI, and say why: a result that cannot be
% regenerated from a shell is not a measurement. scripts/make_project.tcl
% exists for GUI inspection and is not the shipping flow.

\subsection{Vivado Simulator (xsim)}
\label{sec:xsim}
% Driven by scripts/run_sim.sh and scripts/run_all_sims.sh.

\subsection{PyTorch}
\label{sec:pytorch}
% [paszke2019pytorch]. Training and post-training quantisation only - the
% integer pipeline that the hardware matches is the exporter's, not PyTorch's.

\subsection{C Toolchain}
\label{sec:c-toolchain}
% Two distinct toolchains, and the distinction is worth one sentence each:
% the host gcc that builds the bit-exact reference model and the firmware unit
% tests, and the aarch64 bare-metal toolchain for the PS application.

\subsection{Automation Scripts and Version Control}
\label{sec:automation}
% The point of this subsection is the methodology, not the file list. Say:
% every architectural number is recomputed by analyze_workload.py from the
% manifest rather than hard-coded; every figure in this manuscript comes from
% scripts/thesis_numbers.py for the same reason; the checkpoint is tracked in
% Git LFS and its SHA-256 is carried in the manifest so a golden set can never
% again be silently older than the weights it claims to describe (\DD{019}).

\subsection{Vitis Unified Software Platform}
\label{sec:vitis}
% Where the XSA from Section~\ref{sec:board-impl} goes, and how the PS application is
% built against it.
